\documentclass[11pt,a4paper]{article}
\usepackage[utf8]{inputenc}
\usepackage[T1]{fontenc}
\usepackage[spanish]{babel}
\usepackage{hyperref}
\usepackage{geometry}
\geometry{margin=2.5cm}
\usepackage{parskip}

\title{321 eficiencia sin costo}
\author{Alejandro Valenzuela Rivera\\\small Universidad Autónoma de Nuevo León}
\date{\today}
\begin{document}
\maketitle

\begin{quote}
\textbf{Sesgo declarado.} Este trabajo se sostiene desde el sesgo declarado de su autor, y no finge neutralidad. Lo que afirma se puede tumbar, y ese es el punto: el falsificador está escrito, no escondido.
\end{quote}

\section*{Resumen}
La eficiencia sin costo visible no es una virtud: es un acto cuyo efecto nadie cargó.

\section{Por qué se sostiene}
La regla se sostiene mientras alguien la pueda contestar con lo que Terminator (1984) muestra.

\section{De dónde viene}
Tesis/referencias/index.md, y Terminator (1984).

\section{El caso}
Terminator: eficiencia total, sin costo.

\section{Método}
Lo que no dice cuánto costó no se puede contestar: se admira.

\section{Límites, y qué lo tumba}
Un caso donde el costo se calcule después y sirva.

\section{Lo que queda abierto}
\textquestiondown{}Quién cargó el costo de lo que hizo?

\section*{Cómo se leyó esta tesis}
Esto se escribió en el repositorio man-made-empathy, del que este candado es una pieza: candados/06 del caso/321 eficiencia sin costo.md. La máquina que lo compuso lo hizo desde este mismo archivo, y no agregó nada que no estuviera escrito. Si algo de esto es falso, se dice en el mismo sitio, y se corrige con fecha.

\end{document}
